% TEMPLATE-MILC-v3.tex - plantilla maestra UNICA de las guias MILC v3.
%
% La usan las cuatro familias de guias, cada una con su driver:
%   · Grados 6-11          -> scripts/build-guias-g11.py
%   · Bebras               -> scripts/build-guias-bebras.py
%   · Semillero            -> scripts/build-guias-semillero.py
%   · Territorio interior  -> scripts/build-guias-territorio-interior.py
%
% Antes habia tres copias casi identicas (~400 lineas iguales, 6 distintas):
% cada mejora habia que aplicarla tres veces, y en la practica no se hacia
% (el motor de recursos vivia solo en dos de ellas). Lo que varia entre
% familias esta parametrizado en 6 placeholders que cada driver llena
% (se nombran aqui SIN su sintaxis real: el builder reemplaza por texto plano,
% tambien dentro de los comentarios, y un valor multilinea rompe el preambulo):
%
%   HEADER_IZQ        encabezado izquierdo de pagina
%   HEADER_DER        encabezado derecho de pagina
%   PORTADA_ETIQUETA  rotulo sobre el numero grande (GRADO/MOMENTO/MODULO)
%   PORTADA_SERIE     linea de serie de la portada
%   PORTADA_ID        identificador de la guia en la portada
%   PORTADA_CIRCULO   el circulito de grado (°), o vacio si no aplica
%   PDF_TITULO        titulo en texto plano (sin \textbf ni comillas TeX) para
%                     los metadatos del PDF (/Title); lo produce lib_xelatex.texto_pdf
%
% Composicion (auditoria 2026-09-03): las bandas de estacion piden espacio
% (\Needspace*) y prohiben el salto justo despues (\nopagebreak), asi no quedan
% huerfanas al pie; las cajas largas (softbox, drawbox, infoband, puentebox) son
% `breakable`, asi una caja de media pagina ya no arrastra todo a la siguiente
% dejando la anterior al 40 %. Todo texto sobre mostaza o verde va en milcNegro
% (blanco sobre #E5B400 da 1,93:1 y sobre #58A923 2,95:1; ninguno pasa WCAG AA).
% El resto de claves las documenta content/guias/_SCHEMA.md. El script valida
% que no queden placeholders sin reemplazar antes de escribir el archivo final.

% Etiquetado PDF (latex-lab): /Lang, estructura y texto alternativo de las
% figuras, para lectores de pantalla. Debe ir ANTES de \documentclass.
\DocumentMetadata{lang=es-CO,tagging=on}
\documentclass[11pt,letterpaper]{article}
% headheight=24pt: la cabecera derecha lleva dos lineas (identidad de la guia
% y estacion actual). headsep se acorta en la misma medida para que el bloque
% de texto quede exactamente donde estaba con headheight=12pt.
\usepackage[margin=2.0cm,headheight=24pt,headsep=13pt]{geometry}
\usepackage[table]{xcolor}
\usepackage{fontspec}
\usepackage[spanish]{babel}
\usepackage{tikz}
% Librerías TikZ para los diagramas de las guías (bloque `recursos:`):
%  · babel  → babel en español vuelve activos caracteres como «>», y TikZ los usa
%             en sus opciones (>=stealth, ->). Sin esta librería, cualquier
%             diagrama con flechas revienta con «\language@active@arg> has an
%             extra }». Debe ir ANTES de que se dibuje nada.
%  · positioning        → sintaxis «below left=2pt and 3pt».
%  · decorations.pathreplacing → llaves ({brace}) para señalar rangos.
\usetikzlibrary{babel,positioning,decorations.pathreplacing}
\usepackage{graphicx}
% Circuitos: el trabajo de electrónica se hace hoy con micro:bit y sus sensores
% integrados (MakeCode, sin cableado), así que no se necesitan esquemáticos.
% Si algún día entran montajes con protoboard, basta descomentar esta línea:
% los diagramas .tex/.tikz declarados en `recursos:` ya se incrustan solos.
% \usepackage{circuitikz}
\usepackage[most]{tcolorbox}
\usepackage{tabularx}
\usepackage{array}
\usepackage{fancyhdr}
\usepackage{titlesec}
\usepackage[document]{ragged2e}  % bandera a la derecha en todo el cuerpo
\usepackage{enumitem}
\usepackage{microtype}
\usepackage{etoolbox}
\usepackage{needspace}
% hyperref al final del preambulo: metadatos del PDF (titulo, autor, idioma,
% familia), marcadores por estacion (\pdfbookmark) y enlaces sin recuadro.
\usepackage[hidelinks,
  pdftitle={Cosecha P3 — sustentación del mini-estudio},
  pdfauthor={PhD. Álvaro Cárdenas Orozco},
  pdfsubject={Tecnología e Informática · Grado 9 · Guía 30 · MILC v3},
  pdflang={es-CO},
  bookmarksopen=true,bookmarksnumbered=false]{hyperref}

% Helvetica Neue no trae la flecha U+2192, asi que cada «→» escrito literalmente
% en el YAML se compilaba a NADA, en silencio (462 flechas en 61 guias: rutas del
% tipo «paleta → tipografia → lockup» salian rotas). Mapeamos el caracter a su
% version matematica, que si tiene glifo. Asi el autor puede seguir escribiendo
% «→» comodamente en el YAML.
\usepackage{newunicodechar}
\newunicodechar{→}{\ensuremath{\rightarrow}}
% Mismo problema con los simbolos de marca y vinetas: el «✓» de «una palomita ✓»
% salia como cuadrito vacio en el PDF de todas las guias que lo usaban.
\usepackage{pifont}
\newunicodechar{✓}{\ding{51}}
\newunicodechar{✗}{\ding{55}}
\newunicodechar{✔}{\ding{52}}
\newunicodechar{•}{\textbullet}
\newunicodechar{≥}{\ensuremath{\geq}}
\newunicodechar{≤}{\ensuremath{\leq}}

\setmainfont{Helvetica Neue}
\setsansfont{Helvetica Neue}
\setlength{\parindent}{0pt}
\setlength{\parskip}{6pt}
% Sin particion de palabras: en 6.º-7.º una palabra cortada por guion al final
% de linea («Comuni-cacion») frena la lectura mucho mas que una linea corta.
\hyphenpenalty=10000
\exhyphenpenalty=10000
\pretolerance=2000
\tolerance=3000
\setlength{\emergencystretch}{3em}
\linespread{1.10}
\renewcommand{\arraystretch}{1.25}
\setlist{leftmargin=5mm,itemsep=1mm,topsep=1mm}
\newcolumntype{Y}{>{\RaggedRight\arraybackslash}X}

\definecolor{milcMagenta}{HTML}{E6007E}
\definecolor{milcVino}{HTML}{5A0038}
\definecolor{milcTurquesa}{HTML}{0093A5}
\definecolor{milcVerde}{HTML}{58A923}
\definecolor{milcMostaza}{HTML}{E5B400}
\definecolor{milcOcre}{HTML}{B86600}
\definecolor{milcNegro}{HTML}{241816}
\definecolor{milcGris}{HTML}{F7F7F4}
\definecolor{milcLinea}{HTML}{E8E2DD}
\definecolor{milcGradePrimary}{HTML}{7C3AED}
\definecolor{milcGradeDark}{HTML}{4C1D95}
\definecolor{milcGradeSoft}{HTML}{EDE9FE}
\definecolor{milcGradeAccent}{HTML}{CFE7FF}
\definecolor{milcGradeText}{HTML}{FFFFFF}
\color{milcNegro}

\pagestyle{fancy}
\fancyhf{}
% Cabecera de dos lineas: identidad de la guia arriba y, a la derecha y en
% gris, la estacion en curso (\markright desde \milcestacion). Antes de la
% primera estacion la segunda linea va vacia; el \strut mantiene la altura
% para que la linea de arriba no baile entre paginas.
\lhead{\small Tecnología e Informática · Grado 9\\[-1pt]{\footnotesize\strut}}
\rhead{\small Guía 30 · MILC v3\\[-1pt]{\footnotesize\strut\color{milcNegro!65}\rightmark}}
\cfoot{\small\thepage}
\renewcommand{\headrulewidth}{0pt}

% Sobre mostaza (#E5B400) y verde (#58A923) el texto va en milcNegro; sobre el
% resto de la paleta, en blanco. Se usa dentro de fonttitle/fontupper, donde un
% \color posterior manda sobre coltitle.
\newcommand{\milctintasobre}[1]{%
  \ifboolexpr{test{\ifstrequal{#1}{milcMostaza}} or test{\ifstrequal{#1}{milcVerde}}}%
    {\color{milcNegro}}{\color{white}}}

\newtcolorbox{titlebox}[1]{
  enhanced,colback=#1,colframe=#1,arc=7mm,boxrule=0pt,
  left=7mm,right=7mm,top=3mm,bottom=3mm,
  before skip=8pt,after skip=8pt,
  fontupper=\sffamily\bfseries\Large\color{white}
}

\newtcolorbox{softbox}[3]{
  enhanced,breakable,pad at break=3mm,
  colback=#2,colframe=#1,borderline west={4pt}{0pt}{#1},
  arc=2mm,boxrule=.4pt,left=4mm,right=4mm,top=3mm,bottom=3mm,
  before skip=6pt,after skip=8pt,title={#3},coltitle=white,
  colbacktitle=#1,fonttitle=\sffamily\bfseries\milctintasobre{#1},fontupper=\RaggedRight,
  attach boxed title to top left={xshift=3mm,yshift=-2mm},
  boxed title style={arc=2mm, boxrule=0pt}
}

\newtcolorbox{drawbox}[2]{
  enhanced,breakable,pad at break=4mm,
  colback=white,colframe=#1,arc=4mm,boxrule=1pt,
  left=5mm,right=5mm,top=4mm,bottom=4mm,before skip=8pt,after skip=8pt,
  title={#2},coltitle=white,colbacktitle=#1,fonttitle=\sffamily\bfseries\milctintasobre{#1},
  fontupper=\RaggedRight,
  attach boxed title to top left={xshift=4mm,yshift=-2mm},
  boxed title style={arc=3mm, boxrule=0pt}
}

\newtcolorbox{infoband}[1]{
  enhanced,breakable,pad at break=2mm,colback=#1!8,colframe=#1!50,arc=2mm,boxrule=.5pt,
  left=4mm,right=4mm,top=2mm,bottom=2mm,before skip=4pt,after skip=4pt,
  fontupper=\small\RaggedRight
}

% Puente narrativo entre secciones (contrato editorial Conéctate).
% Da continuidad: "Acabas de X. Ahora vas a Y porque Z."
% Visualmente: banda izquierda en color del grado, texto en itálica pequeña.
\newtcolorbox{puentebox}{
  enhanced,breakable,pad at break=2mm,colback=milcGradeSoft,colframe=milcGradeDark,
  borderline west={3pt}{0pt}{milcGradeDark},
  arc=0mm,boxrule=0pt,left=5mm,right=4mm,top=2.5mm,bottom=2.5mm,
  before skip=4pt,after skip=8pt,
  fontupper=\itshape\small\color{milcNegro}\RaggedRight
}
% La flecha va en modo matematico a proposito: Helvetica Neue NO tiene el glifo
% U+2192, asi que \textrightarrow se compilaba a NADA («Missing character: There
% is no → in font Helvetica Neue») y el puente salia sin flecha. En modo
% matematico la flecha viene de la fuente matematica y si se dibuja.
\newcommand{\puente}[1]{%
  \begin{puentebox}%
    \textcolor{milcGradeDark}{$\rightarrow$}\hspace{2mm}#1%
  \end{puentebox}%
}

\newcommand{\linead}[1]{\noindent\color{milcLinea}\rule{#1}{0.5pt}\color{milcNegro}}
\newcommand{\respuesta}[1]{\par\vspace{2mm}\foreach \n in {1,...,#1}{\noindent\color{milcLinea}\rule{\linewidth}{0.5pt}\par\vspace{5mm}}\color{milcNegro}}
\newcommand{\checkbox}{\textcolor{milcGradeDark}{\fbox{\phantom{x}}}\hspace{5pt}}

% ─── Listas aireadas para las guías socioemocionales ────────────────────
% Componer las verificaciones como «\checkbox texto\\» deja el texto que salta
% de línea debajo del cuadrito y apelmaza el bloque. Estos entornos alinean la
% continuación y airean los ítems. Los usa Territorio Interior; cualquier
% familia puede hacerlo.
\newlist{checklista}{itemize}{1}
\setlist[checklista]{
  label=\textcolor{milcGradeDark}{\fbox{\phantom{x}}},
  leftmargin=9mm, labelsep=3.2mm, itemsep=2.4mm,
  topsep=2.5mm, parsep=0pt, partopsep=0pt, before=\RaggedRight
}
\newlist{pasoslista}{enumerate}{1}
\setlist[pasoslista]{
  label=\textcolor{milcGradeDark}{\sffamily\bfseries\arabic*.},
  leftmargin=8mm, labelsep=3mm, itemsep=2.8mm,
  topsep=1mm, parsep=0pt, partopsep=0pt, before=\RaggedRight
}

\newcommand{\phasecard}[4]{
\begin{tcolorbox}[enhanced,colback=#3,colframe=#2,arc=3mm,boxrule=.5pt,height=3.05cm,valign=center,halign=center,left=1.5mm,right=1.5mm,top=2mm,bottom=2mm]
{\sffamily\bfseries\fontsize{22}{24}\selectfont\textcolor{#2}{#1}}\par
{\sffamily\bfseries\small #4}
\end{tcolorbox}
}

% ── Piezas del rediseno 2026 (legibilidad 6.º-11.º) ───────────────────
% Banda de estacion: reemplaza el titlebox macizo. Titulo a la izquierda,
% dato practico (tiempo/modalidad) a la derecha, en una sola linea.
% Nunca huerfana: pide 9 lineas libres (o salta de pagina rellenando con
% \vfill) y prohibe el salto justo despues de la banda. Nueve y no seis:
% la caja partible que sigue necesita ~80pt (titulo adosado, relleno y dos
% lineas) para arrancar en la misma pagina; con menos, tcolorbox la manda
% entera a la siguiente y la banda se queda sola. El penalty 10000 va
% en `after`, ANTES del espacio posterior: un `after skip` (aunque sea 0pt)
% es un \vskip tras la caja, y ese glue es un punto de corte legal que un
% \nopagebreak escrito despues de \end{tcolorbox} ya no cierra. Cada banda
% deja marcador PDF y actualiza la cabecera derecha.
\newcounter{milcestacion}
\newcommand{\milcestacion}[2]{%
\Needspace*{9\baselineskip}%
\stepcounter{milcestacion}%
\pdfbookmark[1]{#2}{milcestacion\themilcestacion}%
\markright{#2}%
\begin{tcolorbox}[enhanced,colback=#1,colframe=#1,arc=2mm,boxrule=0pt,
  left=5mm,right=5mm,top=2.6mm,bottom=2.6mm,before skip=11pt,
  after={\par\nopagebreak\vskip7pt}]
{\sffamily\bfseries\fontsize{15}{18}\selectfont\milctintasobre{#1}#2}
\end{tcolorbox}}

% Tarjeta de paso numerada: sustituye las tablas «Pilar / Aplicacion», que
% eran formato de consulta docente, no de estudiante. El numero va en un
% circulo sobre el borde para que la secuencia se lea de un vistazo.
\newcommand{\milcpaso}[3]{%
\Needspace{4\baselineskip}%
\begin{tcolorbox}[enhanced,colback=white,colframe=milcLinea,arc=1.5mm,boxrule=.9pt,
  left=14mm,right=4mm,top=3mm,bottom=3mm,before skip=4pt,after skip=4pt,
  fontupper=\RaggedRight,
  overlay={\node[anchor=north west,circle,fill=#2,text=white,inner sep=0pt,
    minimum size=7.4mm,font=\sffamily\bfseries\small]
    at ([xshift=3.6mm,yshift=-3.2mm]frame.north west) {#1};}]
#3
\end{tcolorbox}}

% Casilla marcable de tamano util con lapiz (la anterior era \fbox{\phantom{x}},
% de alto variable segun el texto que la seguia).
\newcommand{\milcmarca}{\framebox[3.8mm]{\rule{0pt}{3.4mm}}}

% Fila marcable de lista suelta (checks de la estacion 1).
\newcommand{\milccheckrow}[1]{%
\par\vspace{1.8mm}\noindent
\makebox[6.4mm][l]{\raisebox{-.55mm}{\textcolor{milcNegro}{\milcmarca}}}%
\parbox[t]{\dimexpr\linewidth-6.4mm\relax}{\RaggedRight #1}\par}

% Celda marcable dentro de la rubrica (tabla de autoevaluacion). Misma
% sangria francesa, pero pensada para vivir dentro de una columna X.
\newcommand{\milcopcion}[1]{%
\makebox[5.2mm][l]{\raisebox{-.55mm}{\textcolor{milcNegro}{\milcmarca}}}%
\parbox[t]{\dimexpr\linewidth-5.2mm\relax}{\RaggedRight #1}}

% ── Recursos visuales de la guía ──────────────────────────────────────
% Declarados en el bloque `recursos:` del YAML y emitidos por el builder.
% \guiaFigura   → imagen rasterizada o PDF (foto, carta celeste, montaje).
% \guiaDiagrama → fuente TikZ/circuitikz (.tex) incrustada como vector:
%                 es la vía para circuitos y esquemáticos editables.
% [#1] = texto alternativo (alt) · #2 = ruta relativa al .tex · #3 = pie
% El alt llega al PDF etiquetado (/Alt) y lo lee el lector de pantalla.
\newcommand{\guiaFigura}[3][]{%
\begin{center}
\includegraphics[width=0.88\linewidth,height=0.38\textheight,keepaspectratio,alt={#1}]{#2}\\[1.5mm]
{\footnotesize\itshape\textcolor{milcNegro!75}{#3}}
\end{center}
}

\newcommand{\guiaDiagrama}[2]{%
\begin{center}
\input{#1}\\[1.5mm]
{\footnotesize\itshape\textcolor{milcNegro!75}{#2}}
\end{center}
}

\begin{document}
\thispagestyle{empty}

% ── Portada ───────────────────────────────────────────────────────────
% Antes: un rectangulo de color a pagina completa. Se veia bien en pantalla,
% pero en la fotocopiadora del colegio se comia el toner y salia gris sucio.
% Ahora la banda ocupa el tercio superior y el resto va en blanco.
\begin{tikzpicture}[remember picture,overlay]
  \path[fill=milcGradePrimary]
    (current page.north west) rectangle ([yshift=-10.6cm]current page.north east);

  \node[anchor=north west,text=milcGradeText] at ([xshift=2.0cm,yshift=-1.55cm]current page.north west)
    {\sffamily\bfseries\fontsize{12}{14}\selectfont GRADO};
  \node[anchor=north west,text=milcGradeText,opacity=.85] at ([xshift=2.0cm,yshift=-2.35cm]current page.north west)
    {\sffamily\bfseries\fontsize{8.5}{10}\selectfont SERIE GUÍAS · TECNOLOGÍA E INFORMÁTICA};
  \node[anchor=north east,text=milcGradeText,opacity=.85,align=right] at ([xshift=-2.0cm,yshift=-1.55cm]current page.north east)
    {\sffamily\bfseries\fontsize{9}{12}\selectfont I.E. Sor María Juliana\\Cartago, Valle del Cauca};

  \node[anchor=west,text=milcGradeText] (gradonum) at ([xshift=1.85cm,yshift=-7.1cm]current page.north west)
    {\sffamily\bfseries\fontsize{112}{112}\selectfont 9};
  % El circulito de grado (°) solo aplica a las guias de grado («11°»); en
  % Bebras y Semillero el numero grande es un momento/modulo, no un grado.
  % Se posiciona relativo al nodo `gradonum`, no en coordenadas absolutas.
  \draw[milcGradeText,line width=1.6pt]
    ([xshift=2.0mm,yshift=-4.5mm]gradonum.north east) circle (.15cm);

  \node[anchor=west,text=milcGradeText,text width=11.4cm,align=left] at ([xshift=1.5cm,yshift=0.5cm]gradonum.east)
    {
     {\sffamily\bfseries\fontsize{9}{11}\selectfont\MakeUppercase{Período 3 · Datos --- del registro al insight}}\\[3mm]
     {\sffamily\bfseries\fontsize{21}{25}\selectfont\setlength{\baselineskip}{27pt}Cosecha P3 ---\\[4pt]sustentación del\\[4pt]mini-estudio}\\[3.5mm]
     {\sffamily\bfseries\fontsize{15}{18}\selectfont Guía 30-9-TIC}
    };
\end{tikzpicture}

\vspace*{9.4cm}
\vfill

{\sffamily\bfseries\fontsize{8}{10}\selectfont\textcolor{milcGradeDark}{LO QUE TE LLEVAS HOY}}

\begin{tcolorbox}[enhanced,colback=white,colframe=milcNegro,arc=2mm,boxrule=1.6pt,
  left=5mm,right=5mm,top=4mm,bottom=4mm,before skip=3pt,after skip=12pt,
  fontupper=\RaggedRight\fontsize{11}{15.5}\selectfont]
Presentación de 6 a 8 diapositivas y sustentación de 10 minutos del mini-estudio, con preguntas del grupo respondidas en vivo
\end{tcolorbox}

\begin{tcolorbox}[enhanced,colback=milcGris,colframe=milcGris,arc=2mm,boxrule=0pt,
  left=5mm,right=5mm,top=3.5mm,bottom=3.5mm,after skip=12pt]
{\sffamily\bfseries\fontsize{8}{10}\selectfont\textcolor{milcNegro!55}{ESTA GUÍA ES DE}}\\[3mm]
\begin{tabularx}{\linewidth}{X p{3.2cm} p{3.2cm}}
\linead{\linewidth} & \linead{3.0cm} & \linead{3.0cm}\\[-1mm]
{\fontsize{8.5}{10}\selectfont\textcolor{milcNegro!55}{Nombre completo}} &
{\fontsize{8.5}{10}\selectfont\textcolor{milcNegro!55}{Grupo}} &
{\fontsize{8.5}{10}\selectfont\textcolor{milcNegro!55}{Fecha}}\\
\end{tabularx}
\end{tcolorbox}

\vfill

\noindent\textcolor{milcLinea}{\rule{\linewidth}{1pt}}\\[2mm]
{\sffamily\bfseries\fontsize{9.5}{12}\selectfont Autor: Dr. Álvaro Cárdenas Orozco}\\[1mm]
{\sffamily\fontsize{8.5}{11}\selectfont\textcolor{milcNegro!55}{Metodología MILC · apertura ancestral · pensamiento computacional · triángulo Dussel–estoicismo–Floridi}}

\newpage

\section*{Apertura: Cosecha P3 --- sustentación del mini-estudio}

\begin{softbox}{milcGradeDark}{milcGradeSoft}{Saber ancestral}
En el pueblo misak, en Guambía (Silvia, Cauca), las decisiones grandes no las toma una sola persona. Se llevan al Nu Nakchak, el espacio de las autoridades junto con los shures y las shuras. Ese modo de decidir dejó documentos con fecha. El Mandato de Vida y Permanencia Misak Misak, en Piendamó, 2005. Y la Misak Ley, proclamada el 12 de agosto de 2007. Lo que quedó escrito ahí no es una opinión del momento. Es una decisión sostenida, con testigos y con años encima. La \textbf{cara de exclusión} es que ese «no» colectivo se sostiene frente a normas externas. Y se sostiene a costa de conflicto con el Estado y con terratenientes. No es un recurso retórico: es una disputa territorial viva.\par\smallskip{\footnotesize\textcolor{milcNegro!70}{\textbf{Fuente.} Pueblo Misak. (2007, 12 de agosto). \emph{Misak Ley}. I Encuentro por la defensa de nuestro Derecho Mayor.}}
\end{softbox}

\begin{softbox}{milcTurquesa}{milcGris}{Saber contemporáneo}
Una sustentación de mini-estudio integra el periodo entero en diez minutos. Sigue una estructura estable. Contexto y pregunta de investigación. Datos usados, con su origen, su limpieza y sus límites. Análisis con las herramientas del periodo. Gráficos honestos, con el eje en cero y la fuente a la vista. Tres hallazgos con su decisión, su responsable y su plazo. Y un cierre con la decisión propuesta. Después vienen las preguntas del grupo. No es una exposición decorativa: es una rendición de cuentas.
\end{softbox}

\begin{softbox}{milcMostaza}{milcGris}{Pregunta-puente}
\textit{¿Qué diferencia hay entre decir que no y sostener un no por escrito, con fecha y con testigos? ¿Y qué le falta a una conclusión de datos que nadie tuvo que defender delante de nadie?}
\end{softbox}

\begin{softbox}{milcVerde}{milcGris}{Saber y saber hacer}
Al terminar podrás: (1) \textbf{analizar} qué piezas de tu trabajo del periodo entran en la sustentación y cuáles se descartan; (2) \textbf{evaluar} tu propio estudio antes de exponerlo, anticipando las preguntas duras que puede recibir; (3) \textbf{crear} y defender una presentación de seis a ocho diapositivas que termine en una decisión propuesta.
\end{softbox}



\small
\begin{tabularx}{\textwidth}{>{\bfseries}p{3.0cm}Y}
\rowcolor{milcGradePrimary!12}Autor & Dr. Álvaro Cárdenas Orozco\\
DBA & Tratamiento y visualización honesta de datos para la toma de decisiones (MEN, T\&I)\\
Referentes & Ética del dato · Visualización honesta · Pensamiento computacional\\
Producto final & Presentación de 6 a 8 diapositivas y sustentación de 10 minutos del mini-estudio, con preguntas del grupo respondidas en vivo\\
\end{tabularx}
\normalsize

\puente{Hace tres meses tu pregunta de investigación era una intuición. Hoy es un estudio con datos, análisis y hallazgos. La sustentación es el cierre que convierte nueve sesiones de oficio en algo que otro puede entender.}

\milcestacion{milcGradeDark}{Ruta MILC: así vamos a aprender}
\begin{tabularx}{\textwidth}{>{\centering\arraybackslash}X>{\centering\arraybackslash}X>{\centering\arraybackslash}X>{\centering\arraybackslash}X}
\phasecard{1}{milcVerde}{milcGris}{Escucha\\[-1mm]\small Observo, escucho y pregunto.} &
\phasecard{2}{milcTurquesa}{milcGris}{Entender\\[-1mm]\small Sistematizo y ordeno.} &
\phasecard{3}{milcMagenta}{milcGris}{Hacer\\[-1mm]\small Construyo y aplico.} &
\phasecard{4}{milcVino}{milcGris}{Cierre\\[-1mm]\small Evalúo y reflexiono.}
\end{tabularx}


\puente{Antes de teorizar, vas a hacer un inventario de cosecha. Todo lo que produjiste el periodo sobre la mesa: tabla, fórmulas, filtros, tabla dinámica, gráficos, hallazgos. Reconocer lo cosechado es la primera mitad de la sustentación.}

\milcestacion{milcVerde}{Estación 1 de 3 · Escucha}

\begin{softbox}{milcVerde}{milcGris}{Escena para escuchar}
\textbf{Actividad 1 · ANALIZA --- Inventario de cosecha} (15 min · individual). (1) Pon sobre la mesa todo lo que produjiste en el periodo: tabla original, tabla limpia, fórmulas, filtros, tabla dinámica, gráficos y reporte de hallazgos. (2) Para cada pieza decide si entra en la sustentación o se queda fuera. (3) Recuerda que diez minutos admiten ocho piezas como máximo, así que vas a tener que cortar. (4) Marca las que entran y las que se quedan fuera. (5) Escribe al lado de cada descarte la razón, y distingue las que descartas por decorativas de las que descartas por cariño.
\end{softbox}

{\sffamily\bfseries\fontsize{8}{10}\selectfont\textcolor{milcNegro!55}{MARCA CUANDO LO HAYAS HECHO}}
\milccheckrow{Hice el inventario completo de lo que produje en el periodo.}
\milccheckrow{Marqué qué piezas entran, con un máximo de ocho, y cuáles quedan fuera.}
\milccheckrow{Escribí la razón de cada descarte y reconocí cuáles descartaba por cariño.}

\begin{infoband}{milcVerde}


\textbf{Lo que va al cuaderno (Actividad 1 · ANALIZA).} \textbf{Título:} «Actividad 1 --- Inventario de cosecha». \textbf{Formato:} tabla de tres columnas (pieza producida / entra o no entra / por qué), con una fila por pieza. \textbf{Extensión:} un tercio de página. \textbf{Sabes que terminaste cuando} puedes justificar cada inclusión y cada descarte sin dudar.
\end{infoband}

\puente{Ya viste tu cosecha. Ahora vas a entender la \textbf{anatomía} de una sustentación: sus momentos, los tiempos de cada uno, qué evitar y cómo responder una pregunta dura sin ponerte a la defensiva.}

\milcestacion{milcTurquesa}{Estación 2 de 3 · Entender}

\begin{softbox}{milcTurquesa}{milcGris}{Lo que vas a entender}
\textbf{Actividad 2 · EVALÚA --- Las tres preguntas duras} (20 min · parejas). (1) Con tu pareja, escriban la anatomía de la sustentación con el tiempo de cada momento. (2) Cada uno le presenta al otro su estudio en dos minutos. (3) Quien escucha escribe las tres preguntas más duras que se le ocurran. (4) Cada uno prepara su respuesta a las tres, en una o dos frases, y marca en cuál la respuesta honesta es «eso no lo miré».
\end{softbox}

\milcpaso{1}{milcTurquesa}{La sustentación se descompone en momentos. La \textbf{portada} con tu nombre y la pregunta de investigación. El \textbf{contexto}, en tres líneas. Los \textbf{datos}, con origen, tamaño y limpieza. El \textbf{análisis}, con una tabla dinámica y un gráfico honesto. Los \textbf{tres hallazgos}, uno por diapositiva. La \textbf{decisión propuesta}. Y las \textbf{preguntas}, que van en vivo y no en diapositiva.}
\milcpaso{2}{milcTurquesa}{Toda diapositiva sigue el patrón «una idea, una imagen, pocas palabras». Si tiene tres ideas, se divide en tres. Si tiene un párrafo, se recorta a viñetas. Si tiene jerga, se traduce. La regla que más sirve: la diapositiva apoya tu voz, no la reemplaza. Quien lee la diapositiva en voz alta ya perdió a la audiencia.}
\milcpaso{3}{milcTurquesa}{Todo se juega en una preparación: anticipa las tres preguntas duras que te pueden hacer. Cómo sabes que tu muestra sirve. Qué pasaría si el dato sucio se contara distinto. Por qué tu decisión y no la contraria. Si no tienes respuesta honesta a las tres, es mejor descubrirlo ensayando.}
\milcpaso{4}{milcTurquesa}{Algoritmo: (1) elige las piezas del inventario; (2) produce las diapositivas con la regla de una idea por diapositiva; (3) ensaya en voz alta con cronómetro; (4) escribe las tres preguntas duras y sus respuestas; (5) sustenta; (6) responde con honestidad, incluido el «eso no lo sé» cuando corresponda; (7) anota qué ajustarías.}

\begin{drawbox}{milcTurquesa}{Anatomía de una sustentación de 10 minutos}
\textbf{Portada} (30 s). \\[1mm] \textbf{Contexto} (1 min). \\[1mm] \textbf{Datos y limpieza} (1,5 min). \\[1mm] \textbf{Análisis: tabla dinámica y gráfico honesto} (2 min). \\[1mm] \textbf{Hallazgos 1, 2 y 3} (1 min cada uno). \\[1mm] \textbf{Decisión y cierre} (1 min). \\[1mm] \textbf{Preguntas del grupo} (1 a 2 min). \\[1mm] \textbf{Total:} 10 minutos y las preguntas.
\end{drawbox}

\begin{softbox}{milcMostaza}{milcGris}{Errores comunes y su corrección}
(1) Leer las diapositivas en voz alta, cuando la audiencia podía leerlas sola. (2) Sobrecargar de texto, que es señal de miedo y no de profundidad. (3) Esconder las limitaciones del estudio, que quien escucha va a ver de todos modos. (4) Mostrar cinco gráficos cuando uno honesto basta. (5) Cerrar con un agradecimiento en vez de cerrar con la decisión propuesta.
\end{softbox}

\begin{infoband}{milcTurquesa}


\textbf{Lo que va al cuaderno (Actividad 2 · EVALÚA).} \textbf{Título:} «Actividad 2 --- Las tres preguntas duras». \textbf{Formato:} la anatomía con sus tiempos, las tres preguntas que te hizo tu pareja y tu respuesta a cada una en una o dos frases. \textbf{Extensión:} media página. \textbf{Sabes que terminaste cuando} sabes cuál de las tres respuestas honestas es «eso no lo miré».
\end{infoband}

\puente{Tienes el método. Toca aplicarlo. Produces las diapositivas, ensayas en voz alta con cronómetro y sustentas ante el grupo respondiendo preguntas.}

\milcestacion{milcMagenta}{Estación 3 de 3 · Hacer}

\begin{softbox}{milcMagenta}{milcGris}{Cómo vas a construir tu producto}
\textbf{Actividad 3 · CREA --- Sustenta el mini-estudio} (40 min · individual). (1) Produce las seis a ocho diapositivas siguiendo la anatomía. (2) Declara las limitaciones del estudio en una diapositiva propia, no escondidas al final. (3) Ensaya una vez en voz alta con cronómetro. (4) Sustenta ante el grupo. (5) Responde al menos dos preguntas con honestidad. (6) Anota la retroalimentación recibida y un ajuste que harías si volvieras a sustentar.
\end{softbox}

\milcpaso{1}{milcMagenta}{Tu sustentación se descompone en cuatro piezas: la producción de las diapositivas, el ensayo cronometrado, la entrega ante el grupo y las preguntas respondidas. Sin el ensayo, las otras tres se resienten.}
\milcpaso{2}{milcMagenta}{Sigue el patrón «audiencia que no vio el proceso». Prepárala como si quien escucha no hubiera estado en ninguna sesión del periodo. Eso te obliga a explicar la pregunta, los datos, el análisis y el hallazgo sin atajos. Si tu sustentación solo se entiende porque ya saben, no es comunicación.}
\milcpaso{3}{milcMagenta}{Cada diapositiva responde una sola pregunta de quien escucha. Qué estudiaste. Con qué datos. Qué encontraste. Qué significa. Qué propones. Si una diapositiva no responde ninguna de esas, sobra.}
\milcpaso{4}{milcMagenta}{Algoritmo: (1) produce las diapositivas; (2) ensaya en voz alta con cronómetro; (3) ten a mano las tres preguntas duras con sus respuestas; (4) sustenta; (5) responde con honestidad; (6) recibe la retroalimentación del docente y de al menos un compañero; (7) anota el ajuste que harías.}

\begin{drawbox}{milcMagenta}{Checklist antes de sustentar}
\checkbox Entre seis y ocho diapositivas, con una idea por diapositiva. \\[1mm] \checkbox Un ensayo en voz alta cronometrado. \\[1mm] \checkbox Las tres preguntas duras escritas con su respuesta. \\[1mm] \checkbox Las limitaciones del estudio en una diapositiva propia. \\[1mm] \checkbox El cierre es la decisión propuesta, con responsable y plazo. \\[1mm] \checkbox Sé en qué punto voy a decir «eso no lo miré».
\end{drawbox}

\begin{softbox}{milcMagenta}{milcGris}{Plantilla de trabajo}
\textbf{Portada.} \textit{Nombre y pregunta de investigación.} \\[1mm] \textbf{Contexto.} \textit{Por qué importa, en tres líneas.} \\[1mm] \textbf{Datos.} \textit{Origen, tamaño, limpieza y límites.} \\[1mm] \textbf{Análisis.} \textit{Una tabla dinámica y un gráfico honesto.} \\[1mm] \textbf{Hallazgos.} \textit{Tres, uno por diapositiva.} \\[1mm] \textbf{Cierre.} \textit{Decisión propuesta, responsable y plazo.}
\end{softbox}

\begin{infoband}{milcMagenta}


\textbf{Lo que va al cuaderno (Actividad 3 · CREA).} \textbf{Título:} «Actividad 3 --- Sustentación del periodo 3». \textbf{Formato:} el índice de diapositivas, las tres preguntas anticipadas, la retroalimentación recibida y el ajuste que harías. \textbf{Extensión:} media página. \textbf{Sabes que terminaste cuando} distingues qué sales a mejorar de qué sales a celebrar de tu propio trabajo.
\end{infoband}

\puente{Lo que sigue resume \emph{qué} entregas hoy: la presentación, la sustentación de diez minutos y las preguntas respondidas. El criterio principal es que alguien que no vio tu proceso entienda tu decisión al final de los diez minutos.}

\milcestacion{milcMostaza}{Producto de la sesión}
\textbf{Presentación de 6 a 8 diapositivas, sustentación de 10 minutos y preguntas respondidas}.
\begin{softbox}{milcMostaza}{milcGris}{Criterios de calidad}
(1) Entre seis y ocho diapositivas, con una sola idea en cada una. (2) Tiempo total dentro del rango acordado. (3) Tres hallazgos con su decisión, su responsable y su plazo. (4) Limitaciones del estudio declaradas en una diapositiva propia. (5) Al menos dos preguntas del grupo respondidas con honestidad. (6) Cierre con la decisión propuesta y no con un agradecimiento.
\end{softbox}

\puente{Antes de cerrar, mira la sustentación desde las cinco dimensiones humanas. Sustentar bien es respeto por quien escucha. Sustentar mal desperdicia el oficio de todo un periodo.}

\milcestacion{milcVino}{Evaluación liberadora · cinco dimensiones}

\begin{softbox}{milcVino}{milcGris}{Sentido de la evaluación}
Evaluar no es solo revisar una nota. Es reconocer qué aprendí, qué cambió en mí, qué emociones manejé, cómo me relaciono con la ciudad y la comunidad, y qué le dejaré a quien viene detrás.
\end{softbox}

\textbf{Mi reflexión liberadora en cinco dimensiones.} Completa cada frase iniciada:

\small
\begin{tabularx}{\textwidth}{>{\bfseries\RaggedRight\arraybackslash}p{3.4cm}Y>{\RaggedRight\arraybackslash}p{4.6cm}}
\rowcolor{milcVino}\color{white}Dimensión & \color{white}\bfseries Mi respuesta (frame starter) & \color{white}\bfseries Algo que descubrí\\
1. Personal & Lo que más cambió en mí fue \linead{6.5cm} & \linead{4.2cm}\\
2. Emocional & La emoción más fuerte fue \linead{2.5cm} y la regulé \linead{3cm} & \linead{4.2cm}\\
3. Ciudadana & En democracia esto importa porque \linead{6cm} & \linead{4.2cm}\\
4. Local (Cartago) & En mi barrio sería útil para \linead{6.5cm} & \linead{4.2cm}\\
5. Intergeneracional & A mi abuela/abuelo le diría \linead{6.5cm} & \linead{4.2cm}\\
\end{tabularx}
\normalsize

\begin{infoband}{milcVino}
\textbf{Para la próxima sustentación.} Vuelve a esta tabla en seis meses. Verás que las respuestas cambian — no porque lo aprendido cambie, sino porque tú cambias. La evaluación liberadora es una conversación contigo a través del tiempo.
\end{infoband}


\puente{Y para terminar, tres voces sobre rendir cuentas. Dussel sobre lo que un discurso calla, Marco Aurelio sobre ser en vez de discutir cómo se debería ser, Floridi sobre en qué proyecto humano estamos trabajando.}

\milcestacion{milcGradeDark}{Triángulo de pensamiento · cierre filosófico}

\begin{softbox}{milcVino}{milcGris}{Dussel · lente del nosotros}
\textit{«Hay lengua cotidiana (la de todos los días), lengua de culturas ilustradas, de cultura de masas, lengua de cultura popular… lengua política (que se comprende no por lo que dice sino por lo que calla, contra quién lo dice, cuándo y por qué…).»}

{\footnotesize\itshape\color{milcNegro!70}--- Enrique Dussel · Filosofía de la liberación (1977), §4.2.4.4}

Una sustentación también se comprende por lo que calla. Si no dices de dónde salieron los datos, cuántos eran ni qué descartaste al limpiarlos, quien escucha lo nota igual. Declarar las limitaciones no debilita el estudio: es lo que lo separa de un discurso que se sostiene por lo que oculta.

\textbf{¿Qué estoy dejando fuera de mi sustentación porque no me conviene decirlo?}
\end{softbox}
\linead{\linewidth}\\[1mm]
\linead{\linewidth}

\begin{softbox}{milcMostaza}{milcGris}{Estoicismo · Marco Aurelio · Meditaciones X, 16 (c. 175 d.C.) · cuidado interior}
\textit{«De hoy más, déjate absolutamente de disputar cuál conviene que sea un hombre bueno, sino procura ser tal en realidad.»}

{\footnotesize\itshape\color{milcNegro!70}--- Marco Aurelio · Meditaciones X, 16 (c. 175 d.C.)}

No gastes los diez minutos explicando cómo debería ser un buen estudio de datos. Muestra el tuyo, con lo que tiene y lo que le falta. Y cuando llegue la pregunta que no puedes responder, «eso no lo miré» es más sólido que cualquier rodeo bien armado.

\textbf{¿Dediqué más tiempo a decir qué es un buen análisis que a mostrar el mío?}
\end{softbox}
\linead{\linewidth}\\[1mm]
\linead{\linewidth}

\begin{softbox}{milcTurquesa}{milcGris}{Floridi · lente de la infoesfera}
\textit{«Una de las preguntas políticas apremiantes que enfrentamos en las sociedades de la información avanzadas es: ¿en qué clase de proyecto humano estamos trabajando? (trad. propia)»}

{\footnotesize\itshape\color{milcNegro!70}--- Luciano Floridi · Commentary on the Onlife Manifesto (2015), § 3.1}

Es la pregunta con la que cierra el año. Aprendiste a limpiar una tabla, a filtrarla, a agruparla y a graficarla. La pregunta que queda no es técnica: para qué. Tu decisión propuesta es tu respuesta, y por eso el cierre no puede ser un agradecimiento.

\textbf{Todo lo que aprendí este año a hacer con datos, ¿al servicio de qué lo quiero poner?}
\end{softbox}
\linead{\linewidth}\\[1mm]
\linead{\linewidth}

\begin{infoband}{milcNegro}
\textbf{Nota para el docente.} Las tres son citas textuales y llevan su referencia debajo. Puedes remitir a la obra en clase.
\end{infoband}

\milcestacion{milcVino}{Mi autoevaluación · marca una casilla por fila}

{\small
\setlength{\extrarowheight}{2.2pt}
\begin{tabularx}{\textwidth}{>{\RaggedRight\arraybackslash\bfseries}p{3.0cm}YYY}
\rowcolor{milcVino}\color{white}\bfseries Criterio & \color{white}\bfseries Lo logré & \color{white}\bfseries En proceso & \color{white}\bfseries Necesito apoyo\\[.6mm]
Comprensión del tema & \milcopcion{Explico las ideas con mis palabras.} & \milcopcion{Reconozco las ideas, falta precisión.} & \milcopcion{Necesito repasar lo básico.}\\[.8mm]
\rowcolor{milcGris}Pensamiento computacional & \milcopcion{Usé los pilares al armar mi producto.} & \milcopcion{Usé algunos pilares.} & \milcopcion{Necesito releer la estación 2.}\\[.8mm]
Aplicación y producto & \milcopcion{Mi producto cumple los criterios.} & \milcopcion{Está armado, con ajustes pendientes.} & \milcopcion{Aún está en construcción.}\\[.8mm]
\rowcolor{milcGris}Reflexión 5 dimensiones & \milcopcion{Respondí cada una con honestidad.} & \milcopcion{Respondí algunas con detalle.} & \milcopcion{Mis respuestas son muy generales.}\\[.8mm]
Triángulo de pensamiento & \milcopcion{Conecté las 3 voces con mi trabajo.} & \milcopcion{Conecté al menos una.} & \milcopcion{No las relaciono con mi proceso.}\\[.8mm]
\end{tabularx}}

\vspace{3mm}
\textbf{Hoy entendí que:}\\[1mm]
\linead{\linewidth}\\[1mm]
\linead{\linewidth}

\textbf{Mi compromiso para la próxima sesión es:}\\[1mm]
\linead{\linewidth}\\[1mm]
\linead{\linewidth}

\begin{softbox}{milcMostaza}{milcGris}{Cita de despedida · Marco Aurelio}
\textit{``El obstáculo en el camino se vuelve el camino. Lo que estorba al camino se vuelve camino.''}\\[1mm]
Cada error de hoy, bien observado, se convierte en pieza del camino que pisarás mañana.
\end{softbox}

\small
\begin{tabularx}{\textwidth}{>{\bfseries}p{3.6cm}Y}
\rowcolor{milcGradeDark!12}Nombre completo: & \linead{10cm}\\
Fecha: & \linead{4cm} \quad Curso/grupo: \linead{4cm}\\
Mi firma: & \linead{9cm}\\
\end{tabularx}
\normalsize

\begin{infoband}{milcGradeDark}
\textbf{Para el año entrante.} Conserva esta presentación, tu bitácora y tus tres hallazgos. Piensa en una pregunta cotidiana que merezca un dato: los gastos de la casa, las horas de sueño, la asistencia a una actividad. Con esa, repite el ciclo entero: recoge, limpia, analiza, grafica, escribe el hallazgo y decide. La phronesis se entrena un ciclo a la vez, no en una sustentación final.
\end{infoband}

\end{document}
